
\section{Contexto y Problemática}

En la era digital, estamos constantemente expuestos a un volumen altísimo de información y contenido. Películas, series, videojuegos, música y libros compiten (entre muchas otras) por nuestra atención y tiempo libre, ofreciéndonos innumerables experiencias. Sin embargo, la propia naturaleza de nuestra memoria, la memoria humana, sumada a esta exposición, a menudo dificulta la retención a largo plazo de los detalles específicos y a las emociones vividas durante el consumo de estas obras.

Con mucha más frecuencia de la que nos gustaría, obras que en su momento nos impactaron profundamente se van perdiendo y desvaneciendo en el recuerdo, conservando apenas una impresión general pero olvidando la riqueza de los matices, matices que muchas veces hacen redonda la obra.

Un cinéfilo podría recordar con cariño una película icónica, pero ser incapaz de evocar la emoción exacta que sintió al presencia una escena clave por primera vez, o un fan de los videojuegos podría olvidar los detalles de una misión particularmente memorable. 

Esta pérdida de especificad no solo afecta a la capacidad de rememorar, sino también a la posibilidad de compartir y discutir estas experiencias con la misma intensidad que la primera vez.

El presente trabajo aborda esta problemática, buscando una forma de preservar la viveza y el detalle de estos "recuerdos culturales"

\section{Presentación de LogArt: Tu galería de recuerdos}

Para dar respuesta al desafío de preservar la riqueza de las experiencias culturales, este trabajo de fin de grado propone el diseño y desarrollo de "LogArt: Galería de Recuerdos".

LogArt es una aplicación web concebida y pensada como un repositorio personal y estructurado donde los usuarios pueden documentar y catalogar sus vivencias (y cualquier tipo de anotaciones detalladas) sobre diversas obras artísticas y de entretenimiento.

El objetivo principal de LogArt es ofrecer un espacio digital (en forma de web) donde se puedan registrar no solo datos objetivos -como el minuto exacto de una estrofa de una canción o el nombre de un personaje secundario- sino también e incluso más importante, las impresiones subjetivas, emociones o reflexiones que una obra suscita en un momento determinado.

De esta manera, la aplicación aspira a convertirse en una especie de diario personal de consumo ``cultural'', permitiendo a los usuarios crear un archivo de ``memorias'' sobre las disciplinas que más les apasionan, como pueden ser la literatura, la música o los videojuegos, para posterior consulta y disfrute.

\section{Motivación del Proyecto}

\subsection{Motivación Personal y Origen de la Idea}

La concepción de LogArt surge de una necesidad personal y una anécdota curiosa. La idea que dio origen al proyecto brotó a partir de un comentario realizado por un amigo cercano del autor. En una conversación informal, este amigo mencionó que había comenzado a ver Los Soprano, la serie de televisión favorita del autor. Durante la charla, hizo referencia a una escena concreta de la primera temporada: ``¿Te acuerdas de cuando Tony y Christopher atropellan con el coche al tipo que les debía dinero? Christopher acababa de comprarse el coche, y al atropellar a esa persona se le rompió completamente el parachoques''. Para sorpresa del autor, a pesar considerar dicha serie una obra maestra, apenas conservaba un vago recuerdo del momento descrito.

Esta experiencia desencadenó una reflexión bastante profunda sobre la fragilidad de nuestra memoria y cómo, incluso en obras que nos marcan profundamente, los detalles concretos y las emociones iniciales tienden a diluirse con el paso del tiempo. Se hizo evidente entonces que esta no era una situación aislada, sino una constante en la forma en la que interactuamos con el inmenso universo del entretenimiento.

La idea de crear una herramienta que permitiera ``anclar'' esos momentos, esas sensaciones y esos datos específicos antes de que se olvidaran, se convirtió en la principal fuerza detrás de este proyecto.

\subsection{Relevancia y Potencial Impacto}

Si bien la motivación inicial del proyecto es eminentemente personal, la utilidad de una herramienta como LogArt trasciende el ámbito individual. En un contexto cultural caracterizado por la producción en masa y el consumo acelerado de nuevos contenidos -Videojuegos que se estrenan diariamente\cite{xatakaSaturacionVideojuegos}, temporadas de series que se suceden sin pausa, catálogos musicales en constante expansión-, la capacidad de mantener un recuerdo vívido y detallado de cada experiencia se vuelve un desafío considerable, pero imprescindible.

LogArt nace como una solución relevante para cualquier persona que desee no solo consumir cultura, sino también reflexionar sobre ella y construir un legado personal de sus interacciones con el arte. Al facilitar la organización y revisión de notas, la aplicación puede incluso enriquecer la experiencia del consumidor, fomentar una apreciación más profunda de las obras y sobre todo, servir como una valiosa herramienta de consulta personal.

Se considera que, al igual que el autor encontró una necesidad no cubierta, muchas otras personas podrían beneficiarse de una plataforma que les permita combatir el olvido y preservar la esencia de lo que en algún momento consideraron importante.

\section{Estructura del Documento}

De aquí en adelante, la presente memoria se organiza de la siguiente manera para detallar el desarrollo y los resultados del proyecto LogArt.

El capítulo 2 expone los objetivos concretos que se plantearon al inicio del trabajo.

El capítulo 3 describe las tecnologías herramientas y metodologías de desarrollo empleadas.

El capítulo 4 ofrece una descripción informática exhaustiva y en detalle de la aplicación, incluyendo su arquitectura, diseño e implementación.

Finalmente, el capítulo 5 presenta las conclusiones obtenidas, los objetivos cumplidos, una reflexión sobre el trabajo realizado y las líneas de trabajo futuras.
